\documentclass{article}
\usepackage{../math_template}

\author{Jonathan Kasongo}
\title{Online Math Olympiad 2012}

\begin{document}
\maketitle
\textit{Selected solves from Evan Chen's Online Math Olympiad 2012.}

\begin{problem}{Problem 5}
Two circles have radius 5 and 26. The smaller circle passes through center
of the larger one. What
is the difference between the lengths of the longest and shortest chords of
the larger circle that are
tangent to the smaller circle?
\end{problem}

\begin{solution}{Solution}
The longest chord cannot exceed the diameter of the larger circle which
is $2\times 26 = 52$. This is achievable by taking the diameter of the
larger circle which is tangent to the smaller one.
\vspace{0.2cm}

Let $O$ be the centre of the larger circle. Let $E$ be the foot of the
perpendicular to $AB$ through $O$. Denote $h = OE$, $x = BE$ and $y = AE$.
\begin{center}
\begin{asy}
/* Geogebra to Asymptote conversion, documentation at artofproblemsolving.com/Wiki go to User:Azjps/geogebra */
import graph; size(170);
real labelscalefactor = 0.5; /* changes label-to-point distance */
pen dps = linewidth(0.7) + fontsize(10); defaultpen(dps); /* default pen style */
pen dotstyle = black; /* point style */
real xmin = -27.39094890611301, xmax = 26.96168859328205, ymin = -27.021936890795512, ymax = 27.305623302584824;  /* image dimensions */
pen ffwwqq = rgb(1.,0.4,0.); pen wwzzff = rgb(0.4,0.6,1.);
pair O = (0.,0.), A = (10.,24.), B = (10.,-24.);

filldraw(circle(O, 26.), ffwwqq + opacity(0.05000000074505806),  ffwwqq);
filldraw(circle((5.,0.), 5.), wwzzff + opacity(0.15000000596046448),  wwzzff);
/* draw figures */
draw(O--A);
draw(O--B);
draw(A--B,  linetype("4 4"));
/* dots and labels */
dot(O, dotstyle);
label("$O$", (0.33260533600915065,0.6327031463560191), NE * labelscalefactor);
dot(A, dotstyle);
label("$A$", (10.32322295441296,24.659725683178497), NE * labelscalefactor);
dot(B, dotstyle);
label("$B$", (10.32322295441296,-23.311752302711056), NE * labelscalefactor);
label("$26$", (6.194868566642791,11.531558730069309), NW * labelscalefactor);
label("$26$", (3.71785593398069,-12.49546380675317), NW * labelscalefactor);
dot((10.,0.), dotstyle);
label("$E$", (10.32322295441296,0.6327031463560191), NE * labelscalefactor);
clip((xmin,ymin)--(xmin,ymax)--(xmax,ymax)--(xmax,ymin)--cycle);
/* end of picture */
\end{asy}
\end{center}
Then
$$
26^2 = x^2 + h^2 = y^2 + h^2.
$$
So clearly $x = y$ and $x$ is minimised when $h$ is as large as possible.
Since $OE \perp AB$, point $E$ is actually the point of tangency from the
chord to the smaller circle so $OE$ is a chord of the smaller circle.
Thus $h$ cannot exceed 10, giving $x = y = 24$.
So the final answer is $52 - (24 + 24) = \boxed{4}$. $\blacksquare$
\end{solution}
\newpage

\end{document}
